\begingroup
\emergencystretch=3em
\section{Условие}
\textbf{Цель работы.} Приобретение практических навыков управления потоками
операционной системы и обеспечения синхронизации между ними.

\textbf{Общее задание.} Разработать программу, обрабатывающую данные
в многопоточном режиме с использованием стандартных средств создания
потоков ОС. Максимальное число одновременно работающих потоков задаётся
параметром запуска. Продемонстрировать фактическое количество потоков
средствами операционной системы. Исследовать зависимость ускорения
и эффективности от размера входных данных и количества потоков.

\textbf{Вариант 12.} Выполнить \(K\) последовательных применений фильтра эрозии
и \(K\) последовательных применений фильтра наращивания к матрице
вещественных чисел. Получить две результирующие матрицы.

\textbf{Принятые соглашения реализации.}
\begin{itemize}
\item Используется C++17. Работа с потоками выполняется через интерфейс
\texttt{mt}; в Linux-реализации библиотеки используются функции POSIX Threads.
\texttt{std::thread}, \texttt{std::mutex} и \texttt{std::filesystem} не применяются.
\item Фильтры независимы: обе цепочки преобразований начинаются с исходной
матрицы, а не с результата другого фильтра.
\item Структурирующий элемент --- квадратное окно нечётного размера
\(W\times W\). На границе рассматривается только пересечение окна с матрицей.
\item Входные значения имеют тип \texttt{double} и должны быть конечными.
Размеры матрицы положительны, \(K\geq0\), запрошенное число потоков \(P\geq1\).
\item Ограничение \(P\) включает главный поток. При \(K>0\) число участников
равно \(p=\min(P,N)\), где \(N\) --- число строк. При \(K=0\) дополнительные
потоки не создаются, а обе выходные матрицы равны исходной.
\end{itemize}

\subsection{Входные и выходные данные}
Первая строка входного текстового файла содержит число строк и столбцов
\(N,M\). Далее записываются \(NM\) вещественных значений. Лишние токены,
недостающие элементы, \texttt{NaN} и бесконечности считаются ошибкой.

Программа сохраняет результаты в файлы с суффиксами
\texttt{\_erosion.txt} и \texttt{\_dilation.txt}. Формат каждого результата
совпадает с форматом входа: размеры и элементы матрицы.
\clearpage
\endgroup
